%% frontmatter/abstract.tex
%% One-paragraph abstract, ~180 words, mathematical arc to spectral
%% clustering. Shares its page with the acknowledgements.

\section*{Abstract}
\addcontentsline{toc}{section}{Abstract}

High-dimensional data sets, in which the number of variables is comparable to the number of observations, break the classical guarantees of multivariate statistics: the eigenvalues of a sample covariance matrix no longer converge to their population counterparts, and the corresponding eigenvectors may carry little or no information about the population directions. This report develops the random matrix theory of this regime, from first principles to a working clustering method. We first study pure noise, where the sample eigenvalues spread into the Marchenko-Pastur bulk and the largest eigenvalue exhibits Tracy-Widom fluctuations at the bulk's edge. We then add low-rank signal and obtain the phase transition of the spiked covariance model: a population spike produces an isolated outlier eigenvalue, and its eigenvector aligns with the true direction to an explicit degree, exactly when the spike exceeds the threshold \(1 + \sqrt{y}\) set by the dimension-to-sample-size ratio \(y\). Finally, we show that spectral clustering is governed by the same theory: the similarity matrix of clustered data is itself a spiked model, and the resulting accuracy formula, with no free parameters, matches what we observe on real handwritten digits to within half a percentage point. Monte Carlo simulations test every result along the way.

\vspace{0.45in}
